\section{Correspondence with Regional Parameter Methods}
\label{sec:r8_correspondence}
The appropriate comparison is not between a reward-control method and literature confined to transition probabilities. \citet{quatmann2016} define parametric Markov decision processes and reward properties. \citet{heck2025} develop dependency-sensitive parameter lifting for parametric Markov chains. The latter model class does not include the agent's action choice in the same way, but its treatment of shared parameters is directly relevant. We state the correspondence at the level of the target and information structure.

Write the finite Bellman operator as
\begin{equation}
 (\mathcal T_n^\lambda v)(x)=\max_a\{r_n^\lambda(x,a)
                                    +K_n^\lambda(x,a)v\},\qquad
 (r_n^\lambda,K_n^\lambda)=(1-\lambda)(r_n^0,K_n^0)+\lambda(r_n^1,K_n^1).
 \label{eq:r8_regional_target}
\end{equation}
Augment a state by its date and add a cemetery state carrying the missing mass of a substochastic row. The absorbed state has zero subsequent reward; exit rewards already included in $r_n^\lambda$ must not be paid twice. Set terminal rewards at date $H$. An action-dependent reward can equivalently be assigned to an intermediate action state. These transformations produce a finite-horizon parametric reward MDP. Signed rewards can be accommodated directly in finite-horizon Bellman comparisons; when a nonnegative-reward encoding is desired, a date-specific common shift must also be paid on absorbed paths so that its total is policy independent. These are model correspondences, not a priority claim or an assertion that two implementations are identical.

The distinctions are as follows. In the economic target, $\lambda$ is common across dates and states, and the agent chooses the action before the endpoint-law draw. A rectangular relaxation permits endpoint choices to reset across state--date visits. It therefore loses a dependency that a fixed-$\lambda$ policy must obey. The count-information recursion retains the total number of one endpoint's draws remaining, but does not reveal the next draw before the current action. Its maximization is over one common action for the conditional mixture, not separate actions chosen after observing which endpoint will occur. The extra count information yields an upper value for \eqref{eq:r8_regional_target}. Generalized lifting and dependency-reducing transformations address related losses of parameter dependence; their pMC results should not be silently identified with this controlled reward recursion.

Our corrected chord is a further compression of the count bound. Its signed cross-operator term is
\begin{equation}
 (K_n^a-K_n^b)(V_{n+1}^b-V_{n+1}^a),
 \label{eq:r8_signed_compression}
\end{equation}
followed by positive backward propagation of a statewise correction. It replaces an interior count triangle by one correction vector per nontrivial date, while preserving an upper bound and the coefficient ordering established in Supplement S.10. On the lower side, actual feasible class policies are evaluated under the common parameter. Their Bernstein coefficients preserve the dependency rather than receiving the upper relaxation's extra information. This is why an upper relaxation and a neural training residual cannot substitute for feasible-policy evaluation.

The contribution claimed here is this particular signed compression, its ordering and total policy-bank error analysis, and its measured use in a common-target economic decision certificate. The comparison does not claim that regional expected-reward control, dependency-sensitive abstraction, or policy reuse first becomes possible with NBO. The count and chord timing comparison below is also not a benchmark against the complete software systems of \citet{quatmann2016} or \citet{heck2025}. It is an attributable comparison of the two mathematically matched upper components with the same economic target and lower evidence.
